\vspace{0.7cm}

Este trabajo presenta el diseño e implementación de VeriVoto, un prototipo
SaaS multi-tenant de votación digital para organizaciones. El sistema cubre la
gestión de procesos, el acceso del votante, la emisión del voto, los recibos,
la auditoría y el escrutinio. La arquitectura implementada integra cinco
servicios propios, Keycloak, PostgreSQL y comunicación REST con patrón
\emph{outbox}. La autenticación del votante se configura por proceso mediante
código por correo o cuenta institucional. La batería registrada contiene 49
pruebas de backend y 10 de frontend; las pruebas y la compilación del frontend
finalizaron correctamente. Cuatro pruebas de integración ejecutaron las
migraciones y comprobaciones en PostgreSQL 18.6, incluida la consulta aislada
de censos mediante el repositorio JPA; la autorización HTTP y otros repositorios
siguen pendientes. La cobertura de líneas del backend varía entre
0,0\% y 13,6\%. La sesión de usabilidad con cinco votantes aceptó H1b:
los cinco completaron el voto y obtuvieron el recibo sin ayuda en menos de tres
minutos e interpretaron correctamente la verificación (evidencia descriptiva, sin
inferencia estadística). La prueba de carga local rechazó H1c: con 500 usuarios
simultáneos, el p95 de la emisión superó los dos segundos en las tres
ejecuciones. La revisión manual de accesibilidad sigue pendiente, por lo que
H1a no se acepta ni se rechaza con la evidencia disponible. El
trabajo describe el prototipo y sus límites, sin atribuirle resistencia a la
coacción ni verificabilidad extremo a extremo.

\vspace{0.7cm}

\section*{Abstract}
This work presents the design and implementation of VeriVoto, a multi-tenant
voting Software as a Service prototype for organizations. It supports election
management, voter access, ballot submission, receipts, audit records and
tallying. The implemented architecture integrates five application services,
Keycloak, PostgreSQL and REST communication with an outbox pattern. Voter
authentication is configured per election using email codes or institutional
accounts. The recorded test suite contains 49 backend tests and 10 frontend
tests; frontend tests and build completed successfully. Four integration tests
ran migrations and checks on PostgreSQL 18.6, including organization-scoped
census queries through the JPA repository; HTTP authorization and other
repositories remain unverified. Backend line coverage ranges from 0.0\% to
13.6\%. A usability session with five voters
supported H1b: all five cast their vote and obtained the receipt without help in
under three minutes and interpreted the verification correctly (descriptive
evidence, without statistical inference). The local load test rejected H1c: with
500 simultaneous users, the 95th percentile of vote submission exceeded two
seconds in all three runs. The manual accessibility review remains pending, so
H1a is neither accepted nor rejected on the available evidence. This report
describes the prototype and its limitations and does not claim coercion
resistance or end-to-end verifiability.
